\documentclass[12pt,a4paper]{article}
\usepackage[utf8]{inputenc}
\usepackage[russian]{babel}
\usepackage{amsmath, amssymb}
\usepackage{graphicx}
\usepackage{float}
\usepackage{geometry}
\usepackage{hyperref}

\geometry{
    left=20mm,
    right=20mm,
    top=20mm,
    bottom=20mm
}

\title{Стохастическая динамика нелинейного осциллятора 2-го порядка:\\ Резонансные эффекты и усиленная шумом стабильность}
\author{Дипломная работа}
\date{\today}

\begin{document}

\maketitle

\section{Введение и теоретическая модель}
Численные физические и биологические системы (джозефсоновские контакты, кольцевые лазеры, маятники) с диффузионными флуктуациями описываются стохастическим уравнением Ланжевена второго порядка. В отличие от сильно диссипативных моделей первого порядка, здесь исследуется полная динамика с инертностью:
\begin{equation}
    \ddot{x} + \alpha \dot{x} + \sin(x) = a + A \sin(\omega t) + \xi(t),
    \label{eq:main}
\end{equation}
где:
\begin{itemize}
    \item $x(t)$ — координата системы (фазовая переменная);
    \item $\alpha$ — коэффициент затухания (вязкого трения), определяющий диссипацию энергии;
    \item $a$ — постоянная составляющая внешней силы, задающая наклон профиля потенциала;
    \item $A, \omega$ — амплитуда и частота гармонического переменного сигнала;
    \item $\xi(t)$ — белый гауссов шум с нулевым математическим ожиданием $\langle \xi(t) \rangle = 0$ и корреляционной функцией $\langle \xi(t)\xi(t+\tau) \rangle = D \delta(\tau)$, где $D$ — интенсивность шума.
\end{itemize}

Эффективный потенциал системы имеет вид:
\begin{equation}
    U(x) = -a x - \cos(x).
\end{equation}

Наличие второй производной (массы) существенно усложняет картину стохастических явлений, добавляя эффекты низкочастотной фильтрации и кинетического пролета барьеров. Цель работы — проанализировать среднее время переключения (Mean Switching Time, MST), соответствующее переходу частицы из одной потенциальной ямы в соседнюю (достижение границы $x_{th} = \pi \approx 3.14$) при различных режимах затухания $\alpha$.

\section{Численный метод Хюна}
Для высокой точности интегрирования осциллирующей компоненты СДУ (\ref{eq:main}) разбивается на систему уравнений и решается стохастическим методом Хюна. Это двухэтапный метод предиктора-корректора, который обеспечивает лучшую сходимость по сравнению с методом Эйлера. Среднее время MST вычисляется статистически путем усреднения по ансамблю из множества параллельных реализаций (траекторий) с шагом интегрирования $dt = 0.01$.

\section{Анализ частотного отклика и влияния инерции}
Первая серия вычислительных экспериментов направлена на изучение зависимости среднего времени переключения (MST) от частоты внешнего сигнала $\omega$ при различных значениях вязкого трения $\alpha$. Интенсивность шума зафиксирована на низком уровне $D=0.05$.

\begin{figure}[H]
    \centering
    \includegraphics[width=0.9\textwidth]{images/freq_gamma_sweep.png}
    \caption{Зависимость MST от частоты $\omega$ при различных $\alpha$. Параметры: $A=1.0$, $D=0.05$.}
    \label{fig:freq_gamma}
\end{figure}

Из Рис. \ref{fig:freq_gamma} видно наличие ярко выраженных минимумов, что свидетельствует о явлении резонансной активации. Для систем первого порядка (аналог $\alpha \to \infty$) MST остается малым вплоть до нулевых частот. Однако наличие инерции создает эффект "среза" высоких частот: частица физически не успевает раскачаться за быстро меняющейся силой. При малом затухании ($\alpha=0.2$) этот срез наступает позже, так как система легче накапливает кинетическую энергию. При сильном затухании ($\alpha=5$) кинетическая энергия гасится мгновенно, и частица "запирается" в яме (MST выходит на искусственную полку ограничения времени моделирования).

\section{Динамика при малых частотах и эффект NES}
В данной серии вычислительных экспериментов детально исследовано поведение среднего времени переключения в зависимости от интенсивности шума $D$. Особое внимание уделено эффекту "усиленной шумом стабильности" (NES), заключающемуся во временнóм задерживании (увеличении MST) перескока под действием шумовых флуктуаций при наличии медленно меняющегося надпорогового сигнала. 

Для получения гладких статистических кривых усреднение проводилось не менее чем по 1000 реализациям для каждой точки.

\subsection{Влияние частоты сигнала на проявление NES}
На Рис. \ref{fig:nes_freq_sweep} представлено сравнение кривых MST для различных частот $\omega$ при оптимальном затухании $\alpha=1.0$. Также для сравнения приведена асимптотическая кривая (штриховая линия) статического перескока, соответствующая абсолютно медленному (нулевой частоте) возмущению: отсутствие переменного сигнала $A=0$ при увеличенном суммарном наклоне постоянной составляющей $a' = a + A$.

\begin{figure}[H]
    \centering
    \includegraphics[width=0.9\textwidth]{images/nes_freq_sweep.png}
    \caption{Зависимость MST от интенсивности шума $D$ при различных частотах внешнего сигнала $\omega$ в расширенном диапазоне $D \in [10^{-4}, 10]$. Параметры: $\alpha=1.0$, $a=0.5$, $A=1.0$. Дополнительная кривая обозначает статический предел без сигнала: $A=0, a=1.5$.}
    \label{fig:nes_freq_sweep}
\end{figure}

Расширение диапазона на порядок вниз и вверх позволяет отчетливее разделить три режима. При сверхмалых шумах ($D \lesssim 10^{-3}$) динамика близка к почти детерминированной: кривые с периодическим сигналом располагаются ниже статической асимптотики и задаются фазовой синхронизацией с внешним полем. В промежуточной области ($D \sim 10^{-2}\text{--}1$) сохраняется характерный NES-"горб" для ненулевых частот, что отражает конкуренцию детерминированного выталкивания и стохастической расфазировки. При больших шумах ($D \gtrsim 1$) все кривые переходят к крамерсовскому режиму и демонстрируют тенденцию к сближению: вклад частоты ослабевает, а доминирующим становится термоактивационный механизм перескока.

Для наглядного выделения крайних областей диапазона дополнительно построены два увеличенных фрагмента Рис. \ref{fig:nes_freq_sweep}.

\begin{figure}[H]
    \centering
    \includegraphics[width=0.9\textwidth]{images/nes_freq_sweep_low_zoom.png}
    \caption{Увеличенный низкошумовой фрагмент зависимости MST при $D \in [10^{-4}, 10^{-2}]$.}
    \label{fig:nes_freq_low_zoom}
\end{figure}

\begin{figure}[H]
    \centering
    \includegraphics[width=0.9\textwidth]{images/nes_freq_sweep_high_zoom.png}
    \caption{Увеличенный высокошумовой фрагмент зависимости MST при $D \in [1, 10]$.}
    \label{fig:nes_freq_high_zoom}
\end{figure}

На Рис. \ref{fig:nes_freq_low_zoom} видно, что в пределе очень малого шума различия между частотами определяются в первую очередь эффективностью фазового захвата. На Рис. \ref{fig:nes_freq_high_zoom} кривые для ненулевых частот заметно сближаются, что подтверждает доминирование шумовой (термоактивационной) составляющей над частотной селективностью внешнего сигнала.

\subsection{Зависимость эффекта NES от коэффициента затухания}
На Рис. \ref{fig:nes_alpha_sweep} зафиксированы параметры внешнего сигнала ($\omega=0.1$, $A=1.0$) и постоянного смещения ($a=0.5$), при этом варьируется параметр вязкости среды (затухания) $\alpha$, определяющий меру влияния инерционных эффектов. В качестве статической асимптотики (нижняя граница) приведена кривая переключения при выключенном сигнале ($A=0$) и компенсированном наклоне $a'=1.5$. 

\begin{figure}[H]
    \centering
    \includegraphics[width=0.9\textwidth]{images/nes_alpha_sweep.png}
    \caption{Зависимость MST от интенсивности шума $D$ при различных значениях затухания $\alpha$. Параметры: $\omega=0.1$, $a=0.5$, $A=1.0$. Дополнительная кривая ($A=0, a=1.5$) задает асимптотику статического перескока в отсутствие барьера.}
    \label{fig:nes_alpha_sweep}
\end{figure}

Анализ графиков показывает, что вязкость $\alpha$ кардинально меняет поведение системы при малом шуме и является ключевым фактором возникновения Noise Enhanced Stability. При слабом затухании ($\alpha=0.1$, синяя кривая) частица, накапливая кинетическую энергию, успевает баллистически пролететь барьер на первой же подходящей фазе внешнего сигнала, поэтому эффект NES (задержка перескока шумом) практически отсутствует — кривая монотонно спадает.

Переход к передемпфированному режиму ($\alpha=5.0$, оранжевая кривая) полностью подавляет кинетические пролеты. Система становится настолько вязкой, что для преодоления барьера она должна длительное время находиться в оптимальной фазе сигнала. Даже слабые шумовые флуктуации ($D \sim 0.01$) успешно "сбивают" ее траекторию, что приводит к ярко выраженному росту времени переключения при увеличении шума (участок роста на графике перед спадом Крамерса). Промежуточный случай ($\alpha=1.0$) демонстрирует классический баланс с образованием плато.

\newpage

\section{Термоактивация при подпороговом сигнале}
Третий эксперимент исследует режим, где амплитуда внешнего сигнала $A=0.2$ недостаточна для преодоления барьера детерминированно (Рис. \ref{fig:subthreshold}).

\begin{figure}[H]
    \centering
    \includegraphics[width=0.9\textwidth]{images/subthreshold.png}
    \caption{Зависимость MST от шума при подпороговом сигнале ($A=0.2$).}
    \label{fig:subthreshold}
\end{figure}

Для $\alpha=1.0$ и $\alpha=3.0$ при слабом шуме частица не может вырваться из ямы (MST находится на уровне таймаута). По мере роста $D$ начинает экспоненциально доминировать вероятность термоактивационного перескока (переход Крамерса), и MST резко падает. Синий график ($\alpha=0.1$) демонстрирует, что в режиме слабого трения даже столь малый импульс $A=0.2$ из-за резонансного накопления энергии (отсутствия потерь) оказывается надпороговым, и частица сразу перелетает барьер, почти не воспринимая барьерное сопротивление. 

\newpage

\section{Влияние шума на частотный отклик}
В четвертой серии экспериментов (Рис. \ref{fig:sr_classical}) исследуется частотный отклик при оптимальном затухании $\alpha=1.0$ и различных интенсивностях шума $D$.

\begin{figure}[H]
    \centering
    \includegraphics[width=0.9\textwidth]{images/sr_classical.png}
    \caption{Частотный отклик системы. Зависимость MST от частоты сигнала $\omega$ при различных интенсивностях шума $D$. Параметры: $\alpha=1.0, A=1.0$.}
    \label{fig:sr_classical}
\end{figure}

При сверхмалой интенсивности шума ($D=0.01$) система демонстрирует узкий, глубокий и высокодобротный резонанс: переключение происходит в строгой фазовой синхронизации с внешним сигналом. Однако по мере роста температуры (увеличения $D$) тепловые флуктуации начинают "смывать" потенциальный ландшафт. Детерминированная фазовая когерентность нарушается. При $D=1.0$ (красная кривая) резонансный эффект полностью уничтожается температурным перемешиванием: кривая выполаживается и смещается, переходы становятся сугубо случайными.

\newpage
\section{Чисто шумовая активация: переход Крамерса}
Финальный эксперимент (Рис. \ref{fig:kramers}) проведен при полном отсутствии периодического вынуждающего сигнала ($A=0$). Перескоки частицы через потенциальный барьер осуществляются исключительно за счет энергии флуктуаций (термоактивационные переходы).

\begin{figure}[H]
    \centering
    \includegraphics[width=0.9\textwidth]{images/kramers_pure.png}
    \caption{Модель Крамерса. Зависимость MST от интенсивности шума $D$ без внешнего сигнала ($A=0$) при различных коэффициентах трения $\alpha$.}
    \label{fig:kramers}
\end{figure}

Данный график является превосходной численной иллюстрацией закона Аррениуса для переходов Крамерса: среднее время переключения падает экспоненциально (линейно в логарифмическом масштабе) по мере роста $D$. Коэффициент вязкого трения $\alpha$ входит в предэкспоненциальный множитель формулы Крамерса. Графики показывают, что чем более вязкой является среда ($\alpha=5.0$), тем эффективнее она гасит флуктуации, и тем большая интенсивность шума $D$ требуется для преодоления барьера в пределах заданного окна наблюдения. Режим слабого трения ($\alpha=0.5$), напротив, способствует быстрому шумовому пробою.

\newpage
\section{Заключение}
Модель второго порядка демонстрирует богатство нелинейных резонансных и шумовых явлений. Показано, что инерционный член принципиально меняет фильтрующие свойства системы на высоких частотах и способен полностью нивелировать эффект запирания (позволяя совершать баллистические перескоки) при низких значениях коэффициента затухания. Кроме того, численно подтверждены как классическая теория Крамерса для термоактивационных переходов, так и тонкие эффекты разрушения синхронизации при высоких температурах.

\end{document}